\subsection{Executor Roles and State}

SMC is implemented in the \emph{agent executor}, the program that runs the tool-agent loop. 
Given the current task history, the executor prompts the model for the next tool call, executes that call, records the returned result, and repeats this process until the task finishes or the step budget is exhausted. The executor is therefore distinct from both the language model and the tool environment: it alone decides which action–result pairs become part of the task history seen by future model calls.

SMC augments this executor with two models that run concurrently. 
The \textbf{authoritative actor} is the large model whose tool calls define the real execution of the agent.  
The \textbf{speculative drafter} is a smaller and faster model that runs ahead of the actor, trying to speculate short future tool-call sequences from the committed history $H$.
The executor maintains $H$, consisting of the tool calls and returned results that future actor prompts will condition on, and a live task state $E$, which is the state reached by executing the calls in $H$. 
The drafter's calls are executed in an isolated draft state $E^i$, so they produce returned results without immediately changing the live task state.
We write the drafter's current output as a draft chain $Q=(q_1,\ldots,q_D)$, where each $q_i$ contains a drafted tool call and the result of executing it in the draft state. 
A \emph{macro} is a recurring multi-action pattern mined from past executions. 
A \emph{commit} is the executor operation that appends executed tool calls and results pairs to $H$ and advances $E$ consistently with them. 
% We use \textbf{commit} only as this runtime operation: append tool-call and observation pairs to $H$, advance $E$ consistently with those pairs, and restart speculation from the updated history. 
When no macro commit is performed, $H$ grows in the usual one-step way: the executor appends the authoritative actor's tool call and its returned result. 
SMC changes this loop through a single addition, the \emph{macro commit rule} of Section~\ref{sec:commit}, which under specific conditions lets one actor step be followed by several already executed draft steps.
% A macro commit changes this update rule only after the actor confirms the first drafted call. 
% Specifically, when the actor's next tool call matches the first drafted call, we call this matched call the \emph{anchor call}. 
% The anchor call itself is committed using the actor's own returned result; only the subsequent drafted calls and results can be accepted from the draft chain. 
% If these subsequent draft steps pass the online checks, the executor appends them to $H$ and advances $E$ accordingly.

% For a macro step, the first drafted call is not copied; it must be regenerated by the actor.
% Only the already executed drafted tool-call and observation pairs that follow that matching first call can be copied into $H$.

% This convention makes the experimental comparison straightforward.
% The same abstraction supports two comparisons: a runtime that runs draft work but disables macro commit, and an SMC runtime that additionally allows multi-step commits. 
When the macro commit rule is disabled, the same executor still runs the actor and drafter concurrently and still reuses single-step speculative work: whenever the actor proposes the same call as the first drafted step, the drafter's already executed call and result are committed directly. 
We call this configuration the \emph{SA-only baseline} and use it as the matched-runtime comparison in Section~\ref{sec:results}.
% The concrete model choices, serving replicas, and hardware budgets are experimental settings rather than part of the macro rule.
% We next describe how macros are mined before deployment and how the executor applies this commit rule during execution.

% Some deployments also allocate an \textbf{actor replica} for serving isolation and ordinary one-step overlap. The replica runs the same model class as the authoritative actor. It is not a macro proposer and is not part of the macro-selection rule; it appears only in deployment comparisons that use the same hardware budget.

% SMC is therefore a risk-controlled runtime rule rather than a new model-visible action.
% Offline, it keeps only short tool-call sequences that recur in trusted trajectories or in drafter labels, and selects a library under the same runtime used for deployment.
% Online, it requires the first drafted call to match the actor's verified call, copies only following calls that the drafter has already executed in the current run, and enforces a minimum skipped depth so that a fire can remove critical-path work.
% Additional adapter checks reject impossible commits, such as stale environment copies or incompatible tool arguments, but we treat these as conservative implementation constraints rather than as separately claimed contributions.
% The goal is not to prove trajectory-level equivalence to the sequential run. Instead, SMC defines a conservative fired set whose reliability is increased by offline macro evidence, current-run drafter execution, and first-call agreement with the actor.
% We therefore evaluate both latency and paired quality movement rather than treating macro count as a success metric.

\subsection{Macro Mining}

% SMC mines macros offline from traces of successful tasks on a training subset, and retained  macros form the \emph{macro library}.
SMC first extracts candidate multi-action patterns from successful training traces.
It then evaluates these candidates using drafter predictions from sampled training states and keeps only those that satisfy the selection criteria below.
The selected candidates form the runtime \emph{macro library} $\mathcal{M}$.
Each macro records an ordered tool-call sequence together with the argument structure needed to match it, replacing task-specific values such as user IDs, document IDs, or timestamps with slots. 
This normalization lets one macro cover repeated behavior across tasks without requiring identical concrete arguments.

% A macro is a short tool-call sequence with reliability metadata. Given trusted trajectories for a task family, we sample states along those trajectories and ask the drafter to propose the next few tool calls. 

Because a macro is useful only if the drafter can reproduce it at deployment, we label candidates under the same actor--drafter setting used at runtime. 
For each sampled training state, the drafter proposes a short future action chain, which we compare against a reference future trajectory from that state. 
A proposal is labeled correct when its tool-call sequence matches the reference future trajectory after task-specific argument values are normalized as described above.
% the mining matcher. 
Each candidate macro thus carries two kinds of evidence: how often the pattern occurs, and how reliably the drafter reproduces it when the relevant state is encountered.

% Each proposed sequence is labeled against the trusted future trajectory, producing exact- and skeleton-level agreement signals. 
% The miner then deduplicates by tool-call skeleton, computes support and a lower confidence bound (LCB) over observed correctness, and emits candidate libraries indexed by support and reliability thresholds. 
% This mining target is intentionally more conservative than raw frequency: a frequent sequence is useful only if it is reliable enough to skip actor calls later.

For a candidate macro $m$, let $n_m$ be the number of labeled opportunities and $k_m$ be the number of correct drafter proposals. 
We compute a conservative reliability estimate using the lower $\delta$-quantile of a Beta posterior,

\begin{equation}
    \underline{p}_m =
F^{-1}_{\mathrm{Beta}(k_m+1,\; n_m-k_m+1)}(\delta)
\end{equation}
and retain $m$ only when $n_m \ge n_{\min}$ and $\underline{p}_m \ge \tau$. 
The thresholds $n_{\min}$ and $\tau$ are selected on held-out tasks using the same executor configuration as the final evaluation.
% so the selected library reflects not only offline frequency but also whether the drafter can produce the pattern early enough to matter.

% Mining is a pruning step, not a safety proof: a library entry only says a pattern is frequent and predictable enough to be considered. 
These criteria are used only to filter the candidate macros.
Whether a retained macro is committed is determined separately at runtime by the macro commit rule.

\subsection{Macro Commit Rule}
\label{sec:commit}

% At runtime, the drafter continuously extends the draft chain $Q$ from the current $H$ and the executor checks whether a suffix of $H$ together with a prefix of $Q$ matches a macro in $\mathcal{M}$.
At runtime, the drafter continuously extends the draft chain $Q$ from the current $H$. Consider a retained macro $m=(u_1,\ldots,u_K)$. A runtime alignment matches $u_1,\ldots,u_j$ to a suffix of $H$ and the remaining $\ell=K-j$ actions to the draft prefix $q_1,\ldots,q_\ell$. The library stores $m$ once; $j$ and $\ell$ are determined by the state reached at runtime.
An alignment is only a commit candidate; it does not change $H$ by itself. The executor waits for the actor's next tool call. If that call matches $q_1$, we call $q_1$ the \emph{anchor call}. The drafter has already executed $q_1$, so the executor commits it together with the drafter's returned result; only $q_2,\ldots,q_\ell$ are committed without individual actor confirmation. This skips $\ell-1$ actor decisions, so candidate selection requires $\ell-1 \ge L_{\min}$. The floor cannot be enforced from $K$ alone: for the same four-action macro, aligning one versus two actions with $H$ leaves three versus two draft actions and therefore skips two versus one actor decisions.

% A match is only a commit candidate; it does not change $H$ by itself. 
% The executor waits for the actor's next tool call. If that call matches the first drafted call $q_1$, we call $q_1$ the \emph{anchor call}. 
% The drafter has already executed $q_1$, so the executor commits $q_1$ together with the drafter's returned result, without executing the call a second time; the subsequent already executed draft steps $q_2,\ldots,q_r$ then become eligible for macro commit. 
% The candidate is accepted only if it skips at least $L_{\min}$ actor decisions, i.e., $r-1 \ge L_{\min}$, since part of the matched macro may already lie in $H$. 
% On acceptance, the executor appends $q_2,\ldots,q_r$ with their returned results to $H$ and advances $E$ accordingly.

Online checks then discard stale draft work and apply benchmark-specific action safety. In AppWorld, manually specified API-name rules reject irreversible or unknown calls, while forkable mutations require a matching live-state replay. If all checks pass, the executor appends $q_2,\ldots,q_\ell$ and their returned results to $H$ and advances $E$ accordingly.
The drafter is restarted only when the committed history actually diverges from the draft chain, that is, when the actor's call does not match $q_1$. 
In that case the executor executes the actor's call in the live state, waits for its result, commits this step, discards the draft work, and restarts the drafter from the updated history. Whenever the actor confirms $q_1$, the executor instead reuses the drafter's execution: it removes the committed steps from $Q$ ($q_1$ alone if no macro commit, $q_1,\ldots,q_r$ if one does), and the drafter keeps extending the remainder of the draft chain, which the executor
continues to search for macro matches in later iterations. 
% If no retained macro matches or any online check fails, only the anchor step is committed, and the executor continues as the SA-only baseline.
% Algorithm~\ref{alg:smc-runtime} summarizes this rule.
If candidate selection finds no eligible alignment (including the depth floor), or if an online check rejects the selected actions, only the anchor step is committed and the executor continues as the SA-only baseline.
Algorithm~\ref{alg:smc-runtime} summarizes this rule.



\begin{algorithm}[t]
\caption{\textsc{Speculative Macro Commit} Executor}
\label{alg:smc-runtime}
\small
\begin{algorithmic}[1]
\Require macro library $\mathcal{M}$, live task state $E$, task $x$
\Require authoritative actor $A_{\mathrm{actor}}$, speculative drafter $A_{\mathrm{draft}}$
\Require draft depth $D$, minimum skipped length $L_{\min}$
\State $H \gets [\,]$ \Comment{committed task history}
\State launch $A_{\mathrm{draft}}$ to maintain a draft chain $Q=(q_1,\ldots,q_D)$ from $H$
\State $F_{\mathrm{actor}} \gets \Call{SubmitActor}{A_{\mathrm{actor}}, H, E}$
\While{task not finished and budget remains}
    \State $a_{\mathrm{actor}} \gets \Call{Wait}{F_{\mathrm{actor}}}$
    \Comment{actor's proposed tool call (not yet executed)}
    \State $Q \gets \Call{LatestDraftChain}{}$
    \If{\Call{FirstActionMatch}{$a_{\mathrm{actor}}, q_1$}}
        \State \Call{Commit}{$H,E,q_1$}
        \Comment{reuse the drafter's executed call and result}
        \State $r \gets 1$
        \State $c \gets \Call{SelectMacro}{\mathcal{M}, H, Q, L_{\min}}$
        \If{$c \ne \bot$}
            \State $(m,r') \gets c$
            \Comment{$m$ matches $q_1,\ldots,q_{r'}$ and $r'-1 \ge L_{\min}$}
            \If{\Call{OnlineChecks}{$m,H,(q_2,\ldots,q_{r'}),x,E$}}
                \For{$q \in (q_2,\ldots,q_{r'})$}
                    \State \Call{Commit}{$H,E,q$}
                \EndFor
                \State $r \gets r'$
            \EndIf
        \EndIf
        \State \Call{PopDraftPrefix}{$Q, r$}
        \Comment{drafter keeps extending the remainder of $Q$}
    \Else
        \State $o \gets \Call{Execute}{a_{\mathrm{actor}}, E}$
        \Comment{run the actor's call in the live state}
        \State \Call{Commit}{$H,E,(a_{\mathrm{actor}}, o)$}
        \State discard draft work; restart $A_{\mathrm{draft}}$ from $H$
    \EndIf
    \State $F_{\mathrm{actor}} \gets \Call{SubmitActor}{A_{\mathrm{actor}}, H, E}$
\EndWhile
\end{algorithmic}
\end{algorithm}

 
Overall, unlike registered meta-tools, the actor model never emits a macro call: the actor sees the same tool schema as the baseline, and a macro step is a runtime decision over already executed drafter work, not a new action the model must learn. 
% On the other side, SMC is not a lossless speculative-decoding algorithm: after the anchor, the committed draft steps are not regenerated by the actor. 
% SMC is instead a guarded approximation, and for this reason we evaluate both wall-clock latency and paired task quality. rather than treating the number of macro fires as a success metric.

% \subsection{Online Checks and Approximation}

% SMC fires only when three sources of evidence agree.
% First, the sequence belongs to a retained macro from the offline library.
% Second, the drafter has already executed the later steps in the current run.
% Third, the actor independently produces the same first action, which anchors the macro in the real trajectory.

% The online checks enforce conservative admissibility conditions before commit.
% They reject a macro commit when the draft state is stale, when tool arguments are incompatible with the committed history, when the action type is disallowed by the benchmark adapter, or when the match would skip fewer than $L_{\min}$ actor decisions.
% These checks are not presented as independent sources of speedup.
% They define the set of macro commits that the executor is allowed to attempt.

% SMC is therefore not a lossless speculative-decoding algorithm.
% After the anchor, the remaining committed draft steps are not regenerated one by one by the authoritative actor.
% Instead, SMC is a guarded approximation: it commits only draft work supported by offline macro evidence, current-run drafter execution, and actor agreement on the anchor step.
% For this reason, we evaluate both wall-clock latency and paired task quality, rather than treating the number of macro fires as the main success metric.

% \textbf{Why macro commits are not arbitrary shortcuts.}
% A macro commit is allowed only after the offline library, the current drafter run, and the actor agree on the same local path. The library supplies prior evidence that the sequence is common and reliable. The drafter supplies current-run evidence by executing the later calls before they are copied. The actor supplies online verification by choosing the same first call. Additional runtime checks reject stale or incompatible executions before commit. These checks do not make macro commits formally lossless: the later calls are not regenerated by the actor one by one. Instead, they define a high-precision fired set, whose reliability we measure empirically through paired task outcomes and held-out diagnostics.

% \textbf{A probabilistic view of macro firing.}
% Let $Y=1$ denote that a multi-step commit preserves the final task outcome. SMC fires only under the event $F=L\cap D\cap A\cap K$, where $L$ indicates that the sequence belongs to a retained library macro, $D$ that the drafter has already executed the later calls in the current run, $A$ that the actor verifies the first call, and $K$ that the depth requirement and runtime validity checks pass. The library miner estimates a prior reliability for each macro family, while current-run drafter execution and first-call agreement condition this prior on the task state actually encountered online. We do not assume these signals are independent. Instead, we report empirical precision after filtering stages and estimate the quality impact of $\Pr(Y=1\mid F)$ through held-out mechanism panels and paired full-run outcomes.
